\documentclass[12pt]{article}

% House style preamble
\input{.house-style/preamble.tex}

% Numbers for the case sections, generated from data/derived/ (do not edit)
% Generated by scripts/make_tex_numbers.py from data/derived/. Do not edit by hand.
% Every number in sections/case-lttv.tex comes from here.
\newcommand{\lttvBDUAboveBaseline}{below}
\newcommand{\lttvBDUAboveLTTV}{below}
\newcommand{\lttvBDUBaselineTwenty}{9,766}
\newcommand{\lttvBDUImpactTwenty}{858}
\newcommand{\lttvBDURevDrop}{554}
\newcommand{\lttvBDURevFifteen}{8,919}
\newcommand{\lttvBDURevNineteen}{8,364}
\newcommand{\lttvBDUShareOfBaselineNineteen}{8.4\%}
\newcommand{\lttvBDUShareOfBaselineSixteen}{2.2\%}
\newcommand{\lttvBDUShareTwenty}{9\%}
\newcommand{\lttvBDUVsBaselineMax}{12.8\%}
\newcommand{\lttvBDUVsBaselineMin}{4.0\%}
\newcommand{\lttvBDUVsLTTVMax}{4.9\%}
\newcommand{\lttvBDUVsLTTVMin}{1.9\%}
\newcommand{\lttvBandBDUHi}{2.2}
\newcommand{\lttvBandBDULo}{0.8}
\newcommand{\lttvBandSpecHi}{2.5}
\newcommand{\lttvBandSpecLo}{0.8}
\newcommand{\lttvByopEighteen}{15\%}
\newcommand{\lttvByopSeventeen}{10\%}
\newcommand{\lttvByopSixteen}{5\%}
\newcommand{\lttvCPEBaseMax}{3,140}
\newcommand{\lttvCPEBaseMin}{3,121}
\newcommand{\lttvCPEBaselineTwenty}{3,155}
\newcommand{\lttvCPEImpactTwenty}{399}
\newcommand{\lttvCPEObsMax}{3,427}
\newcommand{\lttvCPEObsMin}{3,204}
\newcommand{\lttvCPEShareOfBaselineNineteen}{11.6\%}
\newcommand{\lttvCPEShareOfBaselineSixteen}{2.6\%}
\newcommand{\lttvCPEShareStated}{18\%}
\newcommand{\lttvCPEShareTwenty}{12.6\%}
\newcommand{\lttvCitedClosuresMult}{2.5}
\newcommand{\lttvCitedUptakeHigh}{35\%}
\newcommand{\lttvCitedUptakeLow}{10\%}
\newcommand{\lttvClosIndHi}{33\%}
\newcommand{\lttvClosIndLo}{6\%}
\newcommand{\lttvClosVICorusHi}{11\%}
\newcommand{\lttvClosVICorusLo}{6\%}
\newcommand{\lttvClosVIHi}{10\%}
\newcommand{\lttvClosVILo}{8\%}
\newcommand{\lttvCutoffBDU}{3.1\%}
\newcommand{\lttvCutoffSpec}{4.1\%}
\newcommand{\lttvDesignCorrectBDU}{86\%}
\newcommand{\lttvDesignCorrectCPE}{63\%}
\newcommand{\lttvDesignCorrectSpec}{99\%}
\newcommand{\lttvExemptSixteen}{55}
\newcommand{\lttvFTEBDU}{5,010}
\newcommand{\lttvFTEProd}{7,180}
\newcommand{\lttvFTESpec}{2,880}
\newcommand{\lttvFTESpecProd}{10,060}
\newcommand{\lttvFTESpecProdShare}{66\%}
\newcommand{\lttvGDPTotal}{1,411.1}
\newcommand{\lttvHistSDBDU}{1.9\%}
\newcommand{\lttvHistSDCPE}{9.1\%}
\newcommand{\lttvHistSDSpec}{2.2\%}
\newcommand{\lttvJobsDirect}{6,830}
\newcommand{\lttvJobsSpinoff}{8,300}
\newcommand{\lttvJobsTotal}{15,130}
\newcommand{\lttvKBDU}{1.3}
\newcommand{\lttvKBDUHi}{2.1}
\newcommand{\lttvKBDULo}{0.5}
\newcommand{\lttvKSpec}{0.1}
\newcommand{\lttvKSpecHi}{0.5}
\newcommand{\lttvKSpecLo}{\textminus 0.2}
\newcommand{\lttvNCalibrations}{9}
\newcommand{\lttvNInconclusiveBDU}{4}
\newcommand{\lttvNInconclusiveSpec}{3}
\newcommand{\lttvNetflixFcEighteen}{57.0}
\newcommand{\lttvNetflixFcFifteen}{45.5}
\newcommand{\lttvNetflixFcGrowth}{25\%}
\newcommand{\lttvNetflixObsEighteen}{58.5}
\newcommand{\lttvNetflixObsFifteen}{43.4}
\newcommand{\lttvNetflixObsGrowth}{35\%}
\newcommand{\lttvPassImplied}{358}
\newcommand{\lttvPayCanFifteen}{3,014}
\newcommand{\lttvPayCanIncrease}{111}
\newcommand{\lttvPayCanNineteen}{3,125}
\newcommand{\lttvPayShareFifteen}{87.8\%}
\newcommand{\lttvPayShareNineteen}{87.3\%}
\newcommand{\lttvReportCanShare}{86\%}
\newcommand{\lttvReportClosuresIndep}{25\%}
\newcommand{\lttvReportClosuresVI}{10\%}
\newcommand{\lttvReportPassThrough}{75\%}
\newcommand{\lttvSigmaGrowthEighteen}{4.2\%}
\newcommand{\lttvSigmaGrowthPooled}{4.0\%}
\newcommand{\lttvSigmaReadBDU}{1.9\%}
\newcommand{\lttvSigmaReadSpec}{2.2\%}
\newcommand{\lttvSigmaTech}{1.3\%}
\newcommand{\lttvSigmaTechPooled}{4.7\%}
\newcommand{\lttvSigmaTechSixteen}{4.4\%}
\newcommand{\lttvSpecAboveBaseline}{above}
\newcommand{\lttvSpecAboveLTTV}{above}
\newcommand{\lttvSpecBaselineTwenty}{4,147}
\newcommand{\lttvSpecImpactTwenty}{970}
\newcommand{\lttvSpecShareOfBaselineNineteen}{21.3\%}
\newcommand{\lttvSpecShareOfBaselineSixteen}{4.8\%}
\newcommand{\lttvSpecShareTwenty}{23\%}
\newcommand{\lttvSpecVsBaselineMax}{4.6\%}
\newcommand{\lttvSpecVsBaselineMin}{0.8\%}
\newcommand{\lttvSpecVsLTTVMax}{28.1\%}
\newcommand{\lttvSpecVsLTTVMin}{9.9\%}
\newcommand{\lttvSpliceTol}{2\%}
\newcommand{\lttvUptakeAprilN}{66,000}
\newcommand{\lttvUptakeAprilShare}{0.6\%}
\newcommand{\lttvUptakeJuneN}{177,000}
\newcommand{\lttvUptakeJuneShare}{1.6\%}


% Update PDF metadata
\hypersetup{
  pdftitle={Fifteen thousand jobs: checking a forecast put to Parliament on unbundling Canadian television},
  pdfkeywords={forecast evaluation, ex post analysis, economic impact studies, broadcasting regulation, Canada}
}

% Working title for Econ Journal Watch (2026-10-03); see DECISIONS.md and submission/venue-decision-2026-10-03.md.
\title{Fifteen thousand jobs: checking a forecast put to Parliament on unbundling Canadian television}
\author{Brett Reynolds \orcidlink{0000-0003-0073-7195}%
\thanks{Contact: \href{mailto:brett.reynolds@humber.ca}{brett.reynolds@humber.ca}}\\
Humber Polytechnic \& University of Toronto}
\date{\today}

\begin{document}
\maketitle

\begin{abstract}
In April 2016 a House of Commons committee heard that new television rules from the Canadian Radio-television and Telecommunications Commission (CRTC) would cost \lttvJobsTotal{} media jobs by 2020. The figure came from a commissioned economic-impact study that set out its assumptions and annual paths in enough detail to be checked. I fixed the forecast record before collecting any outcomes, ran a design analysis, and compared the study's two scenarios with CRTC data for 2016--2019. Specialty and pay channels' revenue, on which most of the job figure depends, stayed at or above the study's no-reform baseline (\(k=\lttvKSpec\), where 1 is the forecast effect), though the verdict turns inconclusive if the baseline is taken to be badly wrong. Distributors' revenue fell further than forecast. The modelled transfer of distributors' losses to Canadian channels didn't appear: payments to those channels rose by \$\lttvPayCanIncrease{} million where the study's own rule implies a cut of about \$\lttvPassImplied{} million. The testimony stated the forecast more strongly than the study did, and cited a share of programming spending that the study's own figures don't reproduce.
\end{abstract}

\noindent\textbf{Keywords:} forecast evaluation; ex post analysis; economic impact studies; broadcasting regulation; Canada

% AI-use disclosure, on page 1 (house style). Models actually used: Claude Opus 5.5
% (drafting, analysis code); a ChatGPT research session produced the initial
% candidate list (noted in the Data and code section). Brett to confirm.
\aidisclosure{Claude Opus 5.5}

% Introduction for the EJW paper. Numbers from sections/numbers-lttv.tex;
% literature figures quoted from the sources with locators. Draft, 2026-10-03.

\section{Introduction}
\label{sec:intro}

On 12 April 2016 a committee of the House of Commons heard that the Canadian Radio-television and Telecommunications Commission's (CRTC's) new rules for television would, \enquote{by 2020}, cost \enquote{some 15,130 media jobs}, \enquote{all of this as the direct result of Let's Talk TV regulatory changes} \citep[time marks 0900--0905]{commons2016chpc8}. The number came from an economic-impact study by the consultancy Nordicity and Peter Miller, an engineer and communications lawyer, prepared for ACTRA, the Canadian Media Guild, the Directors Guild of Canada, Friends of Canadian Broadcasting and Unifor \citep{nordicity2015television}. The CRTC had told distributors to offer a \$25 entry-level package and to sell discretionary channels singly or in small bundles \citep{crtc2015worldofchoice}. The study modelled what that would do to channel revenue, programming spending, jobs and GDP through 2020.

Outcomes for those years are now public. I found no published check of the forecast, in the scholarly literature, in the trade press, or in the authors' later reports for the CRTC (the searches are listed in the replication files). This paper supplies one.

Economists have long compared what was predicted before a regulation with what happened after it. \textcite[ii]{harrington1999accuracy} found that ex ante estimates of direct costs exceeded actual costs for 12 of 25 environmental and occupational-safety rules and fell short for 6, with errors \enquote{driven by both baseline and compliance issues}. A later set of case studies found \enquote{a slight tendency to overestimate both costs and benefits} \citep[285]{morgenstern2018retrospective}. \textcite{simpson2014overestimate} warns that a preponderance of overestimates need not show bias, since skewed cost distributions could produce one anyway. His review notes that the few studies of industry's own estimates found them above outcomes in most cases \citep[319--320]{simpson2014overestimate}.

The work I've cited concerns mainly regulators' estimates of compliance costs and benefits. The forecasts that interested parties put to legislators are a different object. They predict economy-wide jobs and output through input-output multipliers, they travel through testimony rather than regulatory filings, and I know of no routine for revisiting them. \textcite[4]{manski2011certitude} observes that in policy analysis \enquote{Point predictions are common and expressions of uncertainty are rare}. A forecast delivered to a parliamentary committee as a single number of jobs that \enquote{will be lost} is a clear case.

The closest precedent is \textcite{kane2026sffa}, who checked the enrolment forecasts made by universities and by Harvard's expert during the litigation that ended in the US Supreme Court's decision in \textit{Students for Fair Admissions v.\ Harvard}. Two of his choices carry over. He reports results against more than one baseline, and he separates the expert's technical estimate, stated in units no public statistic matches, from the forecasts the institutions put to the Court.

This paper is built on the same separation. The report's technical claim is a conditional projection: if the reforms were implemented and its assumptions held, revenue, programming spending and jobs would follow given paths. The public argument is what the committee was told would happen. The two can fare differently, and here they do.

I fixed the forecast record before downloading any outcome data, ran a fake-data design analysis to see which quantities could tell a world with the forecast effect from one without it, and wrote down how each result would be read \citep{gelman2014types, gelman2020workflow}. The main outcome measure is a single number, \(k\): the observed shortfall below the report's own no-reform baseline, as a multiple of its forecast impact.

Specialty and pay channels' revenue, the link on which most of the job figure depends, didn't fall: \(k=\lttvKSpec\), with an interval from \lttvKSpecLo{} to \lttvKSpecHi{}. That verdict is fragile. It depends on how wrong the report's baseline is taken to be, and it becomes inconclusive under \lttvNInconclusiveSpec{} of the \lttvNCalibrations{} calibrations I examined. Distributors' revenue fell further than forecast, though with a wide interval (\(k=\lttvKBDU\), from \lttvKBDULo{} to \lttvKBDUHi{}).

The break in the chain doesn't depend on the baseline at all. Applied to distributors' actual revenue fall, the report's own pass-through rule implies a cut of about \$\lttvPassImplied{} million in what they paid Canadian channels; those payments rose by \$\lttvPayCanIncrease{} million. Uptake of the entry-level package, a key input, was \lttvUptakeJuneShare{} of subscribers in mid-2016 against an assumed \lttvByopSixteen{} for that year.

The testimony went beyond the report. It turned an economy-wide estimate, more than half of it spin-off employment in other sectors, into media jobs lost as the direct result of the policy, and it cited a share of programming spending that the report's own figures don't reproduce.

Section~\ref{sec:lttv-forecast} sets out the decision, the forecast and the testimony. Section~\ref{sec:lttv-method} describes the check, section~\ref{sec:lttv-results} the results, and section~\ref{sec:lttv-links} the links in the chain. Section~\ref{sec:discussion} discusses what the case shows and what it doesn't.

% Case 2: Let's Talk TV. Numbers come from sections/numbers-lttv.tex
% (generated by scripts/make_tex_numbers.py); quotations from the files named
% in notes/source-verification.md. Draft, 2026-10-03.

\section{Unbundling Canadian television}
\label{sec:lttv}

\subsection{The decision and the forecast}

In March 2015 the CRTC told television distributors to offer, by March 2016, an entry-level service \enquote{priced at no more than \$25 (not including equipment) per month} \citep[para.~26]{crtc2015worldofchoice}. Every discretionary channel was to be available on its own or in small packages by March 2016, and in both forms by December 2016 \citep[para.~47]{crtc2015worldofchoice}. The policy came out of the \textit{Let's Talk TV} proceeding, at which, by one later account, \enquote{virtually all parties expressed concerns about potential impacts of unbundling} \citep[para.~xxii, p.~5]{nordicity2015television}.

Nine months later Nordicity and Peter Miller published \textit{Canadian Television 2020}, prepared for ACTRA, the Canadian Media Guild, the Directors Guild of Canada, Friends of Canadian Broadcasting and Unifor \citep{nordicity2015television}. It modelled two worlds from 2015 to 2020: a baseline with \enquote{technological and consumer behaviour changes, assuming that the pre-existing broadcasting regulatory regime remains}, and the same world with the CRTC's decisions overlaid \citep[para.~ii, p.~1]{nordicity2015television}. Every headline number is the gap between the two.

By 2020, the report concluded, the decisions were \enquote{likely to result in a loss of \lttvJobsTotal{} FTEs of employment in the Canadian economy}, along with \$\lttvGDPTotal{} million of GDP. Of those jobs, \lttvJobsDirect{} were in broadcasting and production and \lttvJobsSpinoff{} in other sectors \citep[para.~xliv and Table~1, p.~12]{nordicity2015television}. Those figures sit at the end of a chain. Subscribers move to cheaper unbundled packages; distributors lose retail revenue and pass most of the loss to Canadian channels as lower wholesale fees; weaker channels close; spending on Canadian programming falls; and input-output multipliers turn the spending into jobs and GDP.

The report states the assumptions that drive each link. Subscribers taking the entry-level service with a few chosen channels make up \lttvByopSixteen{} of all subscribers in 2016, \lttvByopSeventeen{} in 2017 and \lttvByopEighteen{} from 2018 \citep[Table~18, p.~76]{nordicity2015television}. Distributors pass \enquote{75\% of their retail revenue loss to Canadian programming services} \citep[para.~207, p.~77]{nordicity2015television}. Ten per cent of the large integrated companies' Category A and B channels and a quarter of the independents' close by 2020 \citep[paras.~232--233, p.~85]{nordicity2015television}.

By 2020 the model takes \$\lttvSpecImpactTwenty{} million (\lttvSpecShareTwenty{}) a year out of specialty and pay channels' revenue and \$\lttvBDUImpactTwenty{} million (\lttvBDUShareTwenty{}) out of distributors' revenue, against the baseline \citep[paras.~xxvi--xxvii]{nordicity2015television}.

\subsection{What the Heritage Committee heard}

On 12 April 2016 Ian Morrison of Friends of Canadian Broadcasting, accompanied by Miller, put the report to the Commons Standing Committee on Canadian Heritage. Its key findings, he said, were that \enquote{by 2020, some 15,130 media jobs will be lost}, along with a fall in Canadian programming expenditure that was \enquote{18\% of what now exists}, \enquote{all of this as the direct result of Let's Talk TV regulatory changes} \citep[time marks 0900--0905]{commons2016chpc8}.

Each phrase is stronger than the report. The \lttvJobsTotal{} jobs are economy-wide, and \lttvJobsSpinoff{} of them are spin-off employment in other sectors; the \lttvCPEShareStated{} is a share of the report's projected 2020 baseline, not of current spending; and the report itself splits its total into direct and spin-off impact. The figure also fails to reproduce. The report's own baseline-scenario figures put 2020 Canadian programming expenditure at \$\lttvCPEBaselineTwenty{} million, so its \$\lttvCPEImpactTwenty{} million reduction is \lttvCPEShareTwenty{} of baseline, not the \lttvCPEShareStated{} its text states \citep[para.~239; Figs.~8, 9, 43]{nordicity2015television}. I tested eleven candidate denominators built from the report's components; none gives \lttvCPEShareStated.

Morrison also offered a forecast of his own, one month after the entry-level service launched: \enquote{the best estimates I've seen, Mr.~Waugh, are that about 4\% of Canadians will go for it} \citep[time marks 0920--0925]{commons2016chpc8}.

\subsection{How the forecast was checked}

The prediction record (forecasters, quantities, paths by year, stated assumptions) was committed before any outcome series was downloaded. I then re-read the report's two paths from the data labels on its charts, checking every value against the totals in its text, and ran a design analysis by fake-data simulation before opening any CRTC file \citep{gelman2014types, gelman2020workflow}. It asked, for each forecast quantity, whether the outcome could separate a world in which the report's impact occurred from one in which it didn't, given that the baseline itself could be wrong.

It depended on the quantity. With the baseline's error set at the volatility of the report's own 2010--2014 data, specialty and pay revenue would be classified correctly in about \lttvDesignCorrectSpec{} of simulated worlds and distributors' revenue in about \lttvDesignCorrectBDU{}. Canadian programming expenditure, whose historical volatility was \lttvHistSDCPE{} a year, came out at \lttvDesignCorrectCPE{}, so it was set aside in advance as descriptive.

The outcome measure is a single number, \(k\): the observed shortfall below the report's baseline, 2016--2019, as a multiple of the report's own forecast impact, so that \(k=0\) is no effect and \(k=1\) is the forecast. The baseline's error is modelled as a random walk from 2014. The forecasters' own cited inputs (uptake between 10 and 35\%, closures up to two and a half times the assumed rate) define a band of \(k\) that counts as consistent with their forecast. An outcome below the forecast path was to be reported as the forecast exceeded, never as confirmed.

How wrong the baseline was can't be known from the outcomes, so the plan fixed an external check: the report's forecast of US Netflix subscribers, which Canadian regulation couldn't affect, against Netflix's reported US paid memberships \citep{netflix2019form10k}. The report forecast \lttvNetflixFcEighteen{} million for 2018 and Netflix reported \lttvNetflixObsEighteen{} million, which by the pre-stated rule puts the error at \lttvSigmaTech{} a year, below the historical volatility of both revenue series. The CRTC series reproduce the report's 2012--2014 values within 2\% \citep{crtc2017sfs2016, crtc2021sfs2020}.

\subsection{What happened}

Specialty and pay revenue didn't fall. From 2016 to 2019 it ran between \lttvSpecVsBaselineMin{} and \lttvSpecVsBaselineMax{} \lttvSpecAboveBaseline{} the report's no-reform baseline, and between \lttvSpecVsLTTVMin{} and \lttvSpecVsLTTVMax{} \lttvSpecAboveLTTV{} its forecast path. At the pre-stated reading \(k=\lttvKSpec\) (95\% interval \lttvKSpecLo{} to \lttvKSpecHi{}), against a band of \lttvBandSpecLo{} to \lttvBandSpecHi{} implied by the forecasters' own inputs.

Distributors' revenue fell further than forecast. It ran between \lttvBDUVsLTTVMin{} and \lttvBDUVsLTTVMax{} \lttvBDUAboveLTTV{} the report's forecast path, falling from \$\lttvBDURevFifteen{} million in 2015 to \$\lttvBDURevNineteen{} million in 2019. Here \(k=\lttvKBDU\) (\lttvKBDULo{} to \lttvKBDUHi{}), inside the forecasters' band of \lttvBandBDULo{} to \lttvBandBDUHi{}.

Neither verdict is robust to how wrong the baseline is taken to be. Specialty and pay revenue becomes inconclusive once the baseline's error reaches \lttvCutoffSpec{} a year, and distributors' revenue at \lttvCutoffBDU{}. Other defensible readings of the Netflix comparison cross those lines: the 2016 miss gives \lttvSigmaTechSixteen{} a year and a pooled estimate \lttvSigmaTechPooled{}. Part of the 2016 miss is an inherited offset, since the report's 2015 starting figure (\lttvNetflixFcFifteen{} million) was already above Netflix's count (\lttvNetflixObsFifteen{} million) \citep{netflix2018form10k}.

Measured as growth, though, Netflix outran the forecast, by \lttvNetflixObsGrowth{} against \lttvNetflixFcGrowth{} from 2015 to 2018, which puts the error at \lttvSigmaGrowthEighteen{} a year. Across the \lttvNCalibrations{} calibrations I examined, specialty and pay revenue is inconclusive under \lttvNInconclusiveSpec{} and smaller than the forecasters' inputs imply under the rest; distributors' revenue is inconclusive under \lttvNInconclusiveBDU{}. The estimate of \(k\) doesn't move, and no interval for specialty and pay revenue excludes zero \citep{gelman2014statcrisis, gelmanStern2006difference}.

\subsection{The links in the chain}

The report's own inputs can be checked without a counterfactual, and they show where its chain parted from events. Uptake of the entry-level service was small. The CRTC counted \lttvUptakeAprilN{} subscribers five weeks after launch \citep{crtc2016basicpackage} and \lttvUptakeJuneN{} by 30 June 2016 \citep[para.~12]{crtc2016hearing0907}, \lttvUptakeAprilShare{} and \lttvUptakeJuneShare{} of all subscribers, against the report's \lttvByopSixteen{} for 2016. I found no later count.

The modelled pass-through isn't visible in the aggregate. While distributors' revenue fell, what they paid to Canadian channels rose, from \$\lttvPayCanFifteen{} million in 2015 to \$\lttvPayCanNineteen{} million in 2019, and the Canadian share of all affiliation payments stayed near \lttvPayShareFifteen{} \citep[Table~U-T11]{crtc2019cmrbdu}. The aggregates don't show why. Rising fees on a shrinking subscriber base are consistent with it, and so is the CRTC's wholesale code, which the report itself lists as coming into effect in January 2016 \citep[p.~73, n.~82]{nordicity2015television}.

Closures are the one input that came out close to the report. Between \lttvClosVILo{} and \lttvClosVIHi{} of the integrated companies' Category A and B channels stopped operating by 2019--2020, against the report's 10\%. For the independents the bounds run from \lttvClosIndLo{} to \lttvClosIndHi{}, because channels that left licensing for exemption can't be told apart from channels that closed. Canadian programming expenditure, read only descriptively, stayed at or above the baseline path through 2019.

\subsection{What the case shows}

The report's technical claim held in one place and failed in another. At the pre-stated reading, distributors lost revenue on the scale the model predicted, though the entry-level service drew far fewer subscribers than it assumed, so how much of that loss the decisions caused isn't settled here. The chain then broke at the transfer to Canadian channels, whose revenue stayed at or above what the report expected with no reform at all.

Most of the jobs figure put to Parliament comes after that break. Of the \lttvJobsTotal{} FTEs, \lttvFTESpecProd{} (\lttvFTESpecProdShare{}) are in specialty and pay services and independent production, and depend mainly on Canadian channels' revenue and programming spending; the other \lttvFTEBDU{} are distributors' own, and follow from the revenue loss that did occur \citep[Tables~22--23, pp.~91--92]{nordicity2015television}. The public argument went further: the testimony presented an economy-wide, mostly spin-off estimate as media jobs that would be lost, as the direct result of the policy, and quoted a share of programming spending that the report's own figures don't support.

Neither of the authors' later reports for the CRTC that I read returns to these numbers \citep{miller2022rights, nordicity2022harnessing}. The case resembles the forecasts made during the \textit{SFFA} litigation that \textcite{kane2026sffa} checked: the forecasters were interested parties, the forecast was placed before decision-makers, and the outcomes are in public data.

% Discussion, limitations, data and code. Numbers from sections/numbers-lttv.tex.
% Draft, 2026-10-03.

\section{What the case shows}
\label{sec:discussion}

The report's technical claim held in one place and failed in another. At the pre-stated reading, distributors lost more revenue than the model predicted, though the interval is wide and the entry-level service drew far fewer subscribers than the model assumed, so how much of that loss the CRTC's decisions caused isn't settled here. The chain then broke at the transfer to Canadian channels, whose revenue stayed at or above what the report expected with no reform at all.

The jobs figure put to Parliament rests on that chain. Of the \lttvJobsTotal{} FTEs, \lttvFTESpecProd{} (\lttvFTESpecProdShare{}) are in specialty and pay services and independent production, and depend in the model on Canadian channels' revenue and programming spending; the other \lttvFTEBDU{} are distributors' own \citep[Tables~22--23, pp.~91--92]{nordicity2015television}. The staff counts show employment moving against that chain: channels shed more staff than the forecast's direct loss although their revenue held, and distributors didn't shed staff although their revenue fell. Channel jobs did fall, then, but not by the route the model describes, and the distributor jobs the model tied to the revenue loss didn't go. Spin-off employment has no observable counterpart, and production employment remains to be checked.

The public argument went further than the technical claim. The testimony presented an economy-wide estimate, more than half of it spin-off employment, as media jobs that would be lost, as the direct result of the policy, and it quoted a share of programming spending that the report's own figures don't support. On the way from report to committee, a conditional projection with stated assumptions became a point prediction without them. It's a particular case of the point \textcite[4]{manski2011certitude} makes about policy analysis generally, where point predictions are common and statements of uncertainty rare.

The report deserves credit for what made this check possible. It stated its assumptions link by link, published annual paths for both scenarios, cited the range of outside estimates it chose among, and gave the figures behind its headline. A forecast of a single number would have left nothing to check but the number. The useful lesson for those who commission or hear such forecasts is in that contrast. Intermediate predictions that public data can test, an uptake share or a payment flow in the second year, let a forecast be held to account before its end date, and a statement of how wrong the baseline could be shows in advance how much the headline depends on it.

\subsection{Limits of the check}

The paper checks one forecast, and nothing here shows that stakeholder forecasts in general are biased; \textcite{simpson2014overestimate} explains why even a run of overestimates needn't show that. The verdict on channel revenue rests on how wrong the report's technology-only baseline is taken to be. I model that error as a random walk calibrated in several ways, the simplest model that lets the uncertainty be stated, and a structural break in viewing could exceed it. Entry-level uptake is known only for 2016, closures of independent channels can't be separated from moves to exemption, the aggregates show that the pass-through didn't appear without showing why, and the staff counts are a different measure from the report's modelled employment. Canadian programming expenditure moved too much from year to year, before the reforms, to test the forecast on that measure.

None of these limits touches the two findings that need no counterfactual. Payments to Canadian channels rose while the report's own rule says they should have fallen, and the testimony stated more than the report did.

\section*{Data and code}

Every number in this paper can be reproduced from the repository at \url{https://github.com/BrettRey/prediction-checker}, which holds a script that downloads every public source used, the frozen prediction record, the scripts that read the report's figures and check them against its text, the design analysis, and the outcome analysis, with the CRTC series used (Open Government Licence~-- Canada). The analysis plan and each choice made after the outcome data was opened are recorded with dates in the repository. The case was first suggested in a list of candidate cases produced in a ChatGPT research session; every claim taken from that list was checked against primary sources before use.


% The AI-use disclosure is now on page 1 (see \aidisclosure above).
% Uncomment this section only for human/funder acknowledgements:
% \section*{Acknowledgements}
% Thanks to NAME for COMMENTS. This work was supported by FUNDER (GRANT).

\clearpage
\printbibliography

\end{document}
